% v19: current questions only; stable support identifiers.
\subsection{残る六型と，再開する具体的な入口}
\label{ns19:support-agenda}

現在の高分母 support の必要候補は\emph{六型・百二十六標識}である．
葉付き $G_{5,7},G_{5,9}$ は，それぞれ六十標識を含めて排除された．
再開時は最初に \nolinkurl{00_CONTINUATION_POLICY_ja.txt} を読む．
以下は現在の計画であり，面白さとインパクトを最大化するために柔軟に変更する．

全候補は五頂点を使い，最小次数二以上，二点連結，辺数七以上，$b_1\ge3$ である．
各二価頂点は三角形に属する．一つの core に高分母 ray が高々一本なら全分母一となる．
表の $ij$ は無向辺，$N$ は標識付き個数であり，補 graph では孤立頂点を省略する．
\begin{center}
\footnotesize
\setlength{\tabcolsep}{4pt}
\begin{tabular}{crclcr}
型&$|E|$&次数列&$E$&低辺の補 graph&$N$\\\hline
$G_{5,11}$&7&$(43322)$&$\{01,02,03,04,12,13,24\}$&$P_4$&60\\
$G_{5,13}$&7&$(44222)$&$\{01,02,03,04,12,13,14\}$&$K_3$&10\\
$G_{5,14}$&8&$(43333)$&$K_5-\{14,23\}$&$2K_2$&15\\
$G_{5,15}$&8&$(44332)$&$K_5-\{24,34\}$&$P_3$&30\\
$G_{5,16}$&9&$(44433)$&$K_5-\{34\}$&$K_2$&10\\
$G_{5,17}$&10&$(44444)$&$K_5$&$\varnothing$&1
\end{tabular}
\end{center}
全 $1024$ 標識付き graph の検算は
\nolinkurl{verify_support_agenda_v19.py} と
\nolinkurl{support_agenda_verification_v19.json} にある．
候補の個数は整数接着や幾何的 atlas の存在数ではない．

\paragraph{第一候補：book の低三角形の無上界解析．}
$G_{5,13}$ の二価 cores $2,3,4$ は全て高近傍 $0,1$ を持つ．
低辺 $23,24,34$ の node 条件から $L\in\{1,3,4\}$，有向三角形，
少なくとも一つの負枝が必要となる．固定 core の低三角形完成は有限かつ完全に解ける．
次は有限な node offset ごとに，負枝 core の左右への相手が同時に存在する条件を比較する．
この段階で低三角形自体を排除できれば，四高中心 $0,1$ を先に解く必要はない．
低三角形が実在した場合に，初めて中心二 core の row/pair，$q_{01}$，高辺の node と $D$ を加える．
固定 core 自体を動かす無限集合を処理したと誤認しない．

\paragraph{並行候補：七高辺のもう一型．}
$G_{5,11}$ の二価頂点は $3,4$，それぞれの高三角形は $013,024$ である．
共通低 punctures $(13;4),(24;3)$ の通常 reverse は絶対値四以下である．
二つの二高標準形と二つの三高 matching cores を，例外接着と負枝を含めて接続する．
book の固定低 row と $D$ の二次復元，葉型で使った自由分母の線形化は，
仮定を照合して転用する価値がある．

\paragraph{密な四型と局所理論．}
$G_{5,14}=K_{2,2,1}$ と $G_{5,15},G_{5,16},G_{5,17}$ が残る．
$G_{5,15}$ の共通低 puncture は $(23;4)$ であり，二価 core $4$ の負枝も保持する．
三高分母には matching と低 row/column の二次復元が使えるが，四高分母へは自動的に拡張できない．
三高の全高 $D$ 同時非負性は $\beta\le-2$ で未解決であり，rank 二 Schur 系を除かない．
高分母 support $H$ と素数 support $G_\ell$ を区別する．
dense mod 五の二模型の整数 lifting も残し，少なくとも二辺の五進付値一を保つ．

\paragraph{再利用と，解決済みの型．}
$K_4$，二葉三角形，$G_{5,5},G_{5,7},G_{5,9},G_{5,10},G_{5,12}$ は再探索しない．
整数 partner の点・全偶奇半直線による表示と，正整数変数のままの $D$ 線形系を再利用する．
$G_{5,15},G_{5,16},G_{5,17}$ は $K_4$ を部分 graph に持つが，
support 全体が $K_4$ である場合の排除はそこへ移せない．
整数候補が残れば actual weights，basin areas，integer framings，作用切片へ進む．
幾何的同値・別退化・局所置換へ研究の主題を移すことも，成果の価値に応じて判断する．
